% Part: turing-machines
% Chapter: undecidability
% Section: unsolvability-decision-problem

\documentclass[../../../include/open-logic-section]{subfiles}

\begin{document}

\olfileid{tur}{und}{dec}
\olsection{నిర్ణయ సమస్య}

ఇచ్చిన \tetoken{వాక్యం}{!!{sentence}} చెల్లుబాటవుతుందో లేదో నిర్ణయించడానికి ప్రభావక పద్ధతి
ఉంటే, అప్పుడే మాత్రమే మొదటిస్థాయి తర్కం \emph{నిర్ణయించదగినది} అంటాం.
వాస్తవానికి అలాంటి పద్ధతి లేదు: మొదటిస్థాయి వాక్యాల చెల్లుబాటుతనాన్ని
నిర్ణయించే సమస్య పరిష్కరించలేనిది.

ఈ ముఖ్యమైన నిషేధాత్మక ఫలితాన్ని స్థాపించడానికి, నిర్ణయ సమస్యను ట్యూరింగ్
యంత్రంతో పరిష్కరించలేమని నిరూపిస్తాం. అంటే, మొదటిస్థాయి \tetoken{వాక్యం}{!!{sentence}} గల
టేపుపై ప్రారంభించినప్పుడల్లా చివరకు ఆగి, ఆ \tetoken{వాక్యం}{!!{sentence}} చెల్లుబాటవుతుందా
లేదా అనుసరించి $1$ లేదా~$0$ను అవుట్‌పుట్‌గా ఇచ్చే ట్యూరింగ్ యంత్రం లేదని
చూపుతాం. చర్చ్--ట్యూరింగ్ సిద్ధాంతప్రతిపాదన ప్రకారం, గణనీయమైన ప్రతి ప్రమేయం
ట్యూరింగ్-గణనీయం. కాబట్టి ఈ ``చెల్లుబాటుతన ప్రమేయం'' అసలు ప్రభావకంగా
గణనీయమైతే, అది ట్యూరింగ్-గణనీయం అయ్యేది. అది ట్యూరింగ్-గణనీయం కాకపోతే,
ప్రభావకంగానూ గణనీయం కాదు.

నిర్ణయ సమస్య పరిష్కరించలేనిదని నిరూపించే మన వ్యూహం, ఆగే సమస్యను దానికి
తగ్గించడం. దీని అర్థం కిందిది: ట్యూరింగ్ యంత్రం~$M_e$ ఇన్‌పుట్~$w$పై ఆగితే
$h(e,w)=1$, లేకపోతే $h(e,w)=0$ అయ్యే ప్రమేయం~$h$ ట్యూరింగ్-గణనీయం కాదని
నిరూపించాం. మొదటిస్థాయి వాక్యాల చెల్లుబాటుతనాన్ని నిర్ణయించే ట్యూరింగ్ యంత్రం
ఉంటే, $h$ను గణించే ట్యూరింగ్ యంత్రం కూడా ఉంటుందని చూపుతాం. $h$ను ట్యూరింగ్
యంత్రం గణించలేదు కనుక, చెల్లుబాటుతనాన్ని నిర్ణయించే ట్యూరింగ్ యంత్రం కూడా
ఉండదు.\sourcecorrection{OLTETURUND-005}{ఈ తగ్గింపు వివరణలో ఆంగ్ల మూలం ప్రమేయం~$h(e,w)$నే ``halts with output'' అని యంత్రపు చర్యతో వర్ణించింది. ముందరి నిర్వచనానికి అనుగుణంగా, $M_e$ ఆగడాన్ని బట్టి $h$ విలువ $1$ లేదా $0$ అవుతుందని సరిచేశాం.}

ఈ వ్యూహంలో మొదటి అంచె: ప్రతి ఇన్‌పుట్~$w$, ట్యూరింగ్ యంత్రం~$M$కు, $M$
నిర్దేశ సమితినీ ఇన్‌పుట్~$w$నూ సూచించే \tetoken{ఒక వాక్యం}{!!a{sentence}} $!T(M, w)$ను,
``$M$ చివరకు ఆగుతుంది'' అని వ్యక్తం చేసే \tetoken{ఒక వాక్యం}{!!a{sentence}}~$!E(M, w)$ను కింది
షరతు నెరవేరేలా ప్రభావకంగా వర్ణించగలమని చూపడం:
\begin{quote}
  $\Entails !T(M, w) \lif !E(M,w)$ అయితే, అప్పుడే మాత్రమే $M$ ఇన్‌పుట్~$w$కు
  ఆగుతుంది.
\end{quote}
మన నిరూపణలో అధిక భాగం $!T(M, w)$,~$!E(M, w)$ అనే ఈ వాక్యాలను వర్ణించడంలో,
$M$ ఇన్‌పుట్~$w$పై ఆగితే, అప్పుడే మాత్రమే $!T(M, w) \lif !E(M, w)$
చెల్లుబాటవుతుందని నిర్ధారించడంలో ఉంటుంది.

\end{document}
